\documentclass[10pt,a4paper]{amsart}
\usepackage[margin=19mm]{geometry}
\usepackage{amsmath,amssymb,mathtools}
\usepackage[scr=rsfs,cal=euler]{mathalfa}
\usepackage{xfrac}
\usepackage[unicode,hidelinks,bookmarks=false]{hyperref}
\newcommand{\ie}{i.e.\ }
\newcommand{\cf}{cf.\ }
\newcommand{\resp}{respectively\ }
\newcommand{\Subsection}{\subsection}
\providecommand{\nobreakdash}{\nobreak\hbox{-}}
\setlength{\emergencystretch}{2em}
\multlinegap0pt
\usepackage{amssymb}
\usepackage{mathtools}
\usepackage{float}
\usepackage{dsfont}
\usepackage{algorithm}\usepackage{algpseudocode}
\usepackage{tikz}
\newtheorem{proposition}{Proposition}
\def\Re{\mathbb{R}}
\def\bO{\mathbb{O}}
\def\ones{{\sf 1}}
\def\tp{{\hbox{\rm\tiny T}}}
\title{Dual Filter: A Transformer-like Inference \\ Architecture for Hidden Markov Models}
\author{Heng-Sheng Chang, Prashant G. Mehta}
\date{Source: arXiv:2505.00818v2 (CC BY 4.0)}
\begin{document}
\maketitle

\noindent\emph{Source and licence.} Dual Filter: A Transformer-like Inference \\ Architecture for Hidden Markov Models. Version arXiv:2505.00818v2 (CC BY 4.0), distributed by arXiv under the Creative Commons Attribution 4.0 International licence, http://creativecommons.org/licenses/by/4.0/, as recorded in the arXiv licence metadata for this exact version and shown on the abstract page https://arxiv.org/abs/2505.00818v2. Full licence notice: \texttt{licenses/arxiv-2505.00818-CC-BY-4.0.txt}.\par\smallskip

\noindent\textit{Editorial scope.} Counted unit: the article's proposition prop:linear\_algebra (source0.tex lines 2577--2589). The article's complete original proof of it is delivered with it (source0.tex lines 2590--2619, one \texttt{proof} environment of 30 lines). No mathematics is written, repaired or re-derived: the statement and the proof are the article's own lines, in the article's own notation, and no retained line is substituted or re-flowed.  No further source-local context is needed: every cross-reference of every retained line resolves inside the two delivered spans.  Nothing was omitted from any retained span, so the package carries no omission record.  No retained line contains a citation, so the wrapper carries no bibliography and no bibliography entry.  No retained line contains a figure environment, an image-inclusion command or drawing code, so no image is delivered and the package declares no asset dependency.  Editorial formatting only, in addition: an AMS \texttt{amsart} preamble that restates the article's own declarations and macros used by the retained lines, namely the environments \texttt{proposition} on separate counters, together with the standard packages those lines need.  The retained environments print in this document as Proposition 1 in order of appearance, whereas the article prints the same items with the numbering of its own full document.  The complete native source of the article as distributed in the arXiv e-print for this exact version is bundled unchanged as \texttt{sources/177448/source0.tex}.
\begin{proposition}\label{prop:linear_algebra}
  Let $s:\bO\to \Re$. Then there exists unique $(s,\tilde{s}) \in \Re
  \times \Re^m$ such that the following decomposition holds:
  \begin{equation*}
  s(z) = \bar{s} +\tilde{s}^\tp e(z) ,\quad z\in\bO.
\end{equation*}
Explicitly, 
\begin{equation*}
\bar{s}:= \frac{1}{m+1} \sum_{z\in\bO} s(z),\quad \text{and} \quad 
\tilde{s} (i) = (s(i) - \bar{s}),\quad
i=1,2,\hdots,m.
\end{equation*}
\end{proposition}

\begin{proof}
  We have 
  \begin{equation*}
\left( \bar{s} +\tilde{s}^\tp e(z) \right)   = \begin{cases}
\left( \bar{s} + ( s(z) -\bar{s} )\right) = s(z) , &  z=1,2,\hdots,m,\\[5pt]
   \left( \bar{s} - \tilde{s}^\tp \ones \right) = s(0),
     & z=0,
     \end{cases}
 \end{equation*}
 where the last step follows because
 \begin{equation*}
   \tilde{s}^\tp \ones = \sum_{i=1}^m \tilde{s}(i)  = \sum_{i=1}^m
     \left(s(i) - \bar{s} \right) = - m \bar{s}  + \sum_{i=1}^m s(i),\end{equation*}
and therefore,  \begin{equation*}
\bar{s} - \tilde{s}^\tp \ones = \bar{s} - \left(- m \bar{s} +
  \sum_{i=1}^m s(i) \right) = (m+1) \bar{s} - \sum_{i=1}^m s(i) = s(0).
\end{equation*}
The decomposition is unique because $\bar{s} + \tilde{s}^\tp e(z)
\equiv 0$ implies
\begin{align*}
  \bar{s} + \tilde{s}^\tp e(i) = \bar{s} + \tilde{s}(i) &= 0, \quad i=1,2,\hdots,m, \\
  \bar{s} + \tilde{s}^\tp e(0) = \bar{s} - \sum_{i=1}^m \tilde{s}(i) &= 0.
\end{align*}
Summing the first of these equations over $i$,
\begin{equation*}
m\bar{s} + \sum_{i=1}^m \tilde{s}(i) = (m+1) \bar{s} =0,
\end{equation*}
which then also implies $\tilde{s}(i) = -\bar{s} =0$ for
$i=1,2,\hdots,m$.         
\end{proof}

\end{document}
